\subsection{Uniform Calculation Logic}
\label{subsec:calculation-logic}

%%%%%%%%%%%%%%%%%%%%%%%%%%%%%%%%%%%%%%%%%%%%%%%%%%%%%%%%%%%%%%%%%%%%%%%%%%%%%%%
% 10.4: The identical calculation procedure for ALL 144 components
%%%%%%%%%%%%%%%%%%%%%%%%%%%%%%%%%%%%%%%%%%%%%%%%%%%%%%%%%%%%%%%%%%%%%%%%%%%%%%%

The power of the 144-component architecture lies in its \emph{structural symmetry}:
every component---regardless of category, dimension, time, or valence---follows
the identical calculation logic. This section specifies that logic precisely.

\subsubsection{The Core Insight: One Formula, 144 Applications}

\begin{tcolorbox}[colback=blue!5!white, colframe=blue!75!black, title=The Uniform Calculation Principle]
Every component of the 144-component matrix is calculated using the 
\textbf{same five-step procedure}:

\begin{enumerate}
\item Determine the \textbf{reference point} $r^i_d$
\item Observe the \textbf{outcome} $x^i_d$
\item Calculate \textbf{Gain} or \textbf{Pain} relative to reference
\item Apply \textbf{loss aversion} $\lambda^i_d$
\item Apply \textbf{time discounting} $\delta^i_d$
\end{enumerate}

The categories (INU, KNU, IDN) differ only in \emph{parameter values}---not 
in the calculation procedure itself.
\end{tcolorbox}

%%%%%%%%%%%%%%%%%%%%%%%%%%%%%%%%%%%%%%%%%%%%%%%%%%%%%%%%%%%%%%%%%%%%%%%%%%%%%%%
\subsubsection{Step 1: Reference Point Determination}
%%%%%%%%%%%%%%%%%%%%%%%%%%%%%%%%%%%%%%%%%%%%%%%%%%%%%%%%%%%%%%%%%%%%%%%%%%%%%%%

\paragraph{The Reference Point.}
Every utility calculation begins with a reference point $r^i_d$---the 
benchmark against which outcomes are evaluated.

\begin{definition}[Reference Point]
\label{def:reference-point}
The reference point $r^i_d$ for category $i$ and dimension $d$ is the 
expected or normative level of outcome in that dimension according to 
the logic of that category.
\end{definition}

\paragraph{Reference Point Formation by Category.}

\begin{table}[htbp]
\centering
\caption{Reference Point Formation}
\label{tab:reference-formation}
\renewcommand{\arraystretch}{1.4}
\begin{tabular}{@{}lp{5cm}p{5cm}@{}}
\toprule
\textbf{Category} & \textbf{Reference Point Logic} & \textbf{Example (Financial)} \\
\midrule
INU & Status quo, personal expectations, social comparison & 
``I earned \euro{80k} last year, I expect similar'' \\
\addlinespace
KNU & Group norm, fair distribution, collective standard & 
``Our family should have \euro{60k} to live decently'' \\
\addlinespace
IDN & Identity prescription, role expectation, self-concept & 
``A successful manager like me should earn \euro{100k}'' \\
\bottomrule
\end{tabular}
\end{table}

\paragraph{Formal Specification.}
Reference points are functions of context, history, and category-specific factors:

\begin{align}
r^{INU}_d &= f_{INU}(\text{status quo}, \text{expectations}, \text{peer comparison}, \Psi) \\
r^{KNU}_d &= f_{KNU}(\text{group norm}, \text{fairness standard}, \text{distribution}, \Psi) \\
r^{IDN}_d &= f_{IDN}(\text{identity prescription}, \text{role model}, \text{self-image}, \Psi)
\end{align}

Section~\ref{subsec:reference-points} develops these functions in detail.

%%%%%%%%%%%%%%%%%%%%%%%%%%%%%%%%%%%%%%%%%%%%%%%%%%%%%%%%%%%%%%%%%%%%%%%%%%%%%%%
\subsubsection{Step 2: Outcome Observation}
%%%%%%%%%%%%%%%%%%%%%%%%%%%%%%%%%%%%%%%%%%%%%%%%%%%%%%%%%%%%%%%%%%%%%%%%%%%%%%%

\paragraph{The Outcome.}
The outcome $x^i_d$ is the actual state in dimension $d$, evaluated from 
the perspective of category $i$:

\begin{table}[htbp]
\centering
\caption{Outcome Interpretation by Category}
\label{tab:outcome-interpretation}
\renewcommand{\arraystretch}{1.3}
\begin{tabular}{@{}lll@{}}
\toprule
\textbf{Category} & \textbf{Outcome Refers To} & \textbf{Example} \\
\midrule
INU & Individual's own state & My salary is \euro{70k} \\
KNU & Group's aggregate state & Our family income is \euro{90k \\
IDN & Match with identity standard & I earn \euro{70k} (vs. ``should'' \euro{100k}) \\
\bottomrule
\end{tabular}
\end{table}

\paragraph{Key Distinction.}
For INU and IDN, the outcome is the \emph{individual's} state.
For KNU, the outcome is the \emph{group's} state (possibly including distribution).

\begin{align}
x^{INU}_d &= \text{individual outcome in dimension } d \\
x^{KNU}_d &= \text{group outcome in dimension } d \text{ (+ distribution quality)} \\
x^{IDN}_d &= \text{individual outcome in dimension } d \text{ (same as INU)}
\end{align}

Note: INU and IDN observe the same outcome $x_d$ but compare it to 
different reference points, yielding different Gains/Pains.

%%%%%%%%%%%%%%%%%%%%%%%%%%%%%%%%%%%%%%%%%%%%%%%%%%%%%%%%%%%%%%%%%%%%%%%%%%%%%%%
\subsubsection{Step 3: Gain and Pain Calculation}
%%%%%%%%%%%%%%%%%%%%%%%%%%%%%%%%%%%%%%%%%%%%%%%%%%%%%%%%%%%%%%%%%%%%%%%%%%%%%%%

\paragraph{The Gain-Pain Decomposition.}
Given outcome $x^i_d$ and reference point $r^i_d$, we compute:

\begin{definition}[Gain and Pain Calculation]
\label{def:gain-pain-calc}
\begin{align}
G^i_{d,t} &= \max(0, x^i_{d,t} - r^i_d) \\
P^i_{d,t} &= \max(0, r^i_d - x^i_{d,t})
\end{align}
\end{definition}

\paragraph{Properties.}
\begin{itemize}
\item \textbf{Mutual exclusivity:} For any outcome, either $G > 0$ or $P > 0$, never both
\item \textbf{Non-negativity:} Both $G \geq 0$ and $P \geq 0$
\item \textbf{Reference-dependence:} The same objective outcome can be G or P depending on $r$
\end{itemize}

\paragraph{Example: Same Outcome, Different Reference Points.}
Consider an individual earning \euro{70,000}:

\begin{table}[htbp]
\centering
\caption{Same Outcome, Different Categories}
\label{tab:same-outcome}
\renewcommand{\arraystretch}{1.3}
\begin{tabular}{@{}llcccc@{}}
\toprule
\textbf{Category} & \textbf{Reference Point} & $r^i_F$ & $x^i_F$ & $G^i_F$ & $P^i_F$ \\
\midrule
INU & Status quo (last year) & \euro{65k} & \euro{70k} & \euro{5k} & 0 \\
KNU & Family needs & \euro{60k} & \euro{70k} & \euro{10k} & 0 \\
IDN & ``Successful manager'' & \euro{100k} & \euro{70k} & 0 & \euro{30k} \\
\bottomrule
\end{tabular}
\end{table}

The \emph{same} \euro{70k} income produces:
\begin{itemize}
\item A Gain for INU (above status quo)
\item A larger Gain for KNU (above family needs)
\item A Pain for IDN (below identity standard)
\end{itemize}

%%%%%%%%%%%%%%%%%%%%%%%%%%%%%%%%%%%%%%%%%%%%%%%%%%%%%%%%%%%%%%%%%%%%%%%%%%%%%%%
\subsubsection{Step 4: Loss Aversion Application}
%%%%%%%%%%%%%%%%%%%%%%%%%%%%%%%%%%%%%%%%%%%%%%%%%%%%%%%%%%%%%%%%%%%%%%%%%%%%%%%

\paragraph{The Loss Aversion Coefficient.}
Pains are weighted more heavily than Gains:

\begin{equation}
U^i_{d,t}(\text{raw}) = G^i_{d,t} - \lambda^i_d \cdot P^i_{d,t}
\label{eq:raw-utility}
\end{equation}

where $\lambda^i_d > 1$ is the loss aversion coefficient.

\paragraph{Category-Specific Loss Aversion.}
Loss aversion may vary by category:

\begin{table}[htbp]
\centering
\caption{Loss Aversion by Category}
\label{tab:lambda-by-category}
\renewcommand{\arraystretch}{1.3}
\begin{tabular}{@{}lll@{}}
\toprule
\textbf{Category} & \textbf{Typical $\lambda$} & \textbf{Rationale} \\
\midrule
INU & $\lambda^{INU} \approx 2.0$ & Standard Kahneman-Tversky estimate \\
KNU & $\lambda^{KNU} \approx 1.5$--$2.0$ & Group losses may be shared/buffered \\
IDN & $\lambda^{IDN} \approx 2.0$--$3.0$ & Identity threats particularly painful \\
\bottomrule
\end{tabular}
\end{table}

\paragraph{Dimension-Specific Loss Aversion.}
Loss aversion also varies by dimension (from Table~\ref{tab:loss-aversion}):

\begin{itemize}
\item Social rejection ($\lambda_S \approx 2$--$4$) hurts more than financial loss
\item Health decline ($\lambda_P \approx 2$--$3$) is particularly painful
\item Digital access loss ($\lambda_D \approx 1.5$--$2$) is less severe
\end{itemize}

\paragraph{The Combined Loss Aversion Parameter.}
In practice, we use:
\begin{equation}
\lambda^i_d = \bar{\lambda}_d \cdot \kappa^i
\end{equation}
where $\bar{\lambda}_d$ is the dimension-specific base and $\kappa^i$ is 
a category-specific multiplier.

%%%%%%%%%%%%%%%%%%%%%%%%%%%%%%%%%%%%%%%%%%%%%%%%%%%%%%%%%%%%%%%%%%%%%%%%%%%%%%%
\subsubsection{Step 5: Time Discounting}
%%%%%%%%%%%%%%%%%%%%%%%%%%%%%%%%%%%%%%%%%%%%%%%%%%%%%%%%%%%%%%%%%%%%%%%%%%%%%%%

\paragraph{The Discount Factor.}
Future utility is discounted:

\begin{equation}
U^i_{d,t}(\text{discounted}) = (\delta^i_d)^t \cdot U^i_{d,t}(\text{raw})
\label{eq:discounted-utility}
\end{equation}

where $\delta^i_d \in (0, 1)$ is the discount factor.

\paragraph{Time Horizon Interpretation.}

\begin{table}[htbp]
\centering
\caption{Time Horizons and Typical Discount Factors}
\label{tab:time-discounting}
\renewcommand{\arraystretch}{1.3}
\begin{tabular}{@{}llcc@{}}
\toprule
\textbf{Horizon} & \textbf{Approximate Duration} & \textbf{$t$} & \textbf{$\delta^t$ (if $\delta = 0.9$)} \\
\midrule
Instant & Now & 0 & 1.00 \\
Short-term & Days to weeks & 1 & 0.90 \\
Medium-term & Months to years & 2 & 0.81 \\
Long-term & Years to lifetime & 3 & 0.73 \\
\bottomrule
\end{tabular}
\end{table}

\paragraph{Category-Specific Discounting.}
Discount rates may vary by category:

\begin{table}[htbp]
\centering
\caption{Discounting by Category}
\label{tab:delta-by-category}
\renewcommand{\arraystretch}{1.3}
\begin{tabular}{@{}lll@{}}
\toprule
\textbf{Category} & \textbf{Typical $\delta$} & \textbf{Rationale} \\
\midrule
INU & $\delta^{INU} \approx 0.85$--$0.95$ & Standard personal discounting \\
KNU & $\delta^{KNU} \approx 0.90$--$0.98$ & Group/future generations valued more \\
IDN & $\delta^{IDN} \approx 0.80$--$0.95$ & Identity can be stable or fragile \\
\bottomrule
\end{tabular}
\end{table}

\paragraph{Context-Dependent Discounting.}
From the $\Gamma$-matrix framework, discounting is context-sensitive:
\begin{equation}
\delta^i_d(\Psi) = \bar{\delta}^i_d + \gamma_{\delta,T} \cdot \Psi_T
\end{equation}

Higher temporal stability ($\Psi_T$) increases patience (higher $\delta$).

%%%%%%%%%%%%%%%%%%%%%%%%%%%%%%%%%%%%%%%%%%%%%%%%%%%%%%%%%%%%%%%%%%%%%%%%%%%%%%%
\subsubsection{The Complete Component Formula}
%%%%%%%%%%%%%%%%%%%%%%%%%%%%%%%%%%%%%%%%%%%%%%%%%%%%%%%%%%%%%%%%%%%%%%%%%%%%%%%

\paragraph{Putting It Together.}
Combining all five steps, the utility contribution of any component is:

\begin{equation}
\boxed{
U^i_{d,t} = (\delta^i_d)^t \cdot \left[ G^i_{d,t} - \lambda^i_d \cdot P^i_{d,t} \right]
}
\label{eq:component-formula}
\end{equation}

where:
\begin{align*}
G^i_{d,t} &= \max(0, x^i_{d,t} - r^i_d) \\
P^i_{d,t} &= \max(0, r^i_d - x^i_{d,t})
\end{align*}

\paragraph{The Formula Applies to All 144 Components.}
This single formula, with category- and dimension-specific parameters,
generates all 144 utility components:

\begin{table}[htbp]
\centering
\caption{Formula Application Examples}
\label{tab:formula-examples}
\small
\renewcommand{\arraystretch}{1.2}
\begin{tabular}{@{}lllll@{}}
\toprule
\textbf{Component} & $r^i_d$ & $x^i_d$ & $\lambda^i_d$ & $\delta^i_d$ \\
\midrule
$U^{INU}_{F,0}$ & \euro{80k} (status quo) & \euro{85k} & 2.0 & 1.0 \\
$U^{KNU}_{S,2}$ & 0.7 (group cohesion norm) & 0.5 & 1.8 & 0.85 \\
$U^{IDN}_{P,1}$ & ``Fit'' (BMI 24) & BMI 28 & 2.5 & 0.90 \\
\bottomrule
\end{tabular}
\end{table}

%%%%%%%%%%%%%%%%%%%%%%%%%%%%%%%%%%%%%%%%%%%%%%%%%%%%%%%%%%%%%%%%%%%%%%%%%%%%%%%
\subsubsection{The Calculation Algorithm}
%%%%%%%%%%%%%%%%%%%%%%%%%%%%%%%%%%%%%%%%%%%%%%%%%%%%%%%%%%%%%%%%%%%%%%%%%%%%%%%

\paragraph{Pseudocode.}
For any component $(i, d, t)$:

\begin{verbatim}
FUNCTION calculate_component(i, d, t, x, Psi):
    
    # Step 1: Determine reference point
    r = get_reference_point(i, d, Psi)
    
    # Step 2: Observe outcome (already given as x)
    
    # Step 3: Calculate Gain and Pain
    IF x > r THEN
        G = x - r
        P = 0
    ELSE
        G = 0
        P = r - x
    END IF
    
    # Step 4: Apply loss aversion
    lambda = get_loss_aversion(i, d, Psi)
    U_raw = G - lambda * P
    
    # Step 5: Apply time discounting
    delta = get_discount_factor(i, d, Psi)
    U_discounted = (delta ^ t) * U_raw
    
    RETURN U_discounted
    
END FUNCTION
\end{verbatim}

\paragraph{Vectorized Form.}
For computational efficiency, the entire 144-component vector can be 
computed in parallel:

\begin{equation}
\mathbf{U} = \boldsymbol{\Delta}^t \odot \left( \mathbf{G} - \boldsymbol{\Lambda} \odot \mathbf{P} \right)
\end{equation}

where:
\begin{itemize}
\item $\mathbf{U} \in \mathbb{R}^{144}$: utility vector
\item $\mathbf{G}, \mathbf{P} \in \mathbb{R}^{144}$: gain and pain vectors
\item $\boldsymbol{\Lambda} \in \mathbb{R}^{144}$: loss aversion vector
\item $\boldsymbol{\Delta}^t \in \mathbb{R}^{144}$: discount factors raised to power $t$
\item $\odot$: element-wise multiplication
\end{itemize}

%%%%%%%%%%%%%%%%%%%%%%%%%%%%%%%%%%%%%%%%%%%%%%%%%%%%%%%%%%%%%%%%%%%%%%%%%%%%%%%
\subsubsection{Parameter Summary}
%%%%%%%%%%%%%%%%%%%%%%%%%%%%%%%%%%%%%%%%%%%%%%%%%%%%%%%%%%%%%%%%%%%%%%%%%%%%%%%

\begin{table}[htbp]
\centering
\caption{Summary of Calculation Parameters}
\label{tab:parameter-summary}
\renewcommand{\arraystretch}{1.4}
\begin{tabular}{@{}llll@{}}
\toprule
\textbf{Parameter} & \textbf{Symbol} & \textbf{Varies By} & \textbf{Typical Range} \\
\midrule
Reference point & $r^i_d$ & Category, Dimension, Context & Domain-specific \\
Loss aversion & $\lambda^i_d$ & Category, Dimension, Context & 1.5--4.0 \\
Discount factor & $\delta^i_d$ & Category, Dimension, Context & 0.80--0.98 \\
\bottomrule
\end{tabular}
\end{table}

\paragraph{Context Modulation.}
All parameters are modulated by context $\Psi$:
\begin{align}
r^i_d(\Psi) &= \bar{r}^i_d + \sum_j \gamma^r_{dj} \cdot \Psi_j \\
\lambda^i_d(\Psi) &= \bar{\lambda}^i_d - \gamma_{\lambda,I} \cdot \Psi_I \\
\delta^i_d(\Psi) &= \bar{\delta}^i_d + \gamma_{\delta,T} \cdot \Psi_T
\end{align}

This context-dependence is developed fully in Section~\ref{subsec:context-modulation}.

%%%%%%%%%%%%%%%%%%%%%%%%%%%%%%%%%%%%%%%%%%%%%%%%%%%%%%%%%%%%%%%%%%%%%%%%%%%%%%%
\subsubsection{Worked Example: Complete Calculation}
%%%%%%%%%%%%%%%%%%%%%%%%%%%%%%%%%%%%%%%%%%%%%%%%%%%%%%%%%%%%%%%%%%%%%%%%%%%%%%%

\paragraph{Scenario.}
A middle manager (salary \euro{70k}) evaluates a job offer at \euro{85k} 
that would require relocating away from family.

\paragraph{Selected Components.}

\textbf{Component: $U^{INU}_{F,2}$ (Individual Financial, Medium-term)}
\begin{align*}
r^{INU}_F &= \text{\euro{70k}} \quad \text{(status quo)} \\
x^{INU}_{F,2} &= \text{\euro{85k}} \quad \text{(new offer)} \\
G^{INU}_{F,2} &= \text{\euro{85k}} - \text{\euro{70k}} = \text{\euro{15k}} \\
P^{INU}_{F,2} &= 0 \\
\lambda^{INU}_F &= 2.0, \quad \delta^{INU}_F = 0.9 \\
U^{INU}_{F,2} &= (0.9)^2 \cdot [\text{\euro{15k}} - 2.0 \cdot 0] = 0.81 \cdot \text{\euro{15k}} = \text{\euro{12.15k}}
\end{align*}

\textbf{Component: $U^{KNU}_{S,1}$ (Collective Social, Short-term)}
\begin{align*}
r^{KNU}_S &= 0.8 \quad \text{(current family cohesion)} \\
x^{KNU}_{S,1} &= 0.5 \quad \text{(after relocation)} \\
G^{KNU}_{S,1} &= 0 \\
P^{KNU}_{S,1} &= 0.8 - 0.5 = 0.3 \\
\lambda^{KNU}_S &= 2.5, \quad \delta^{KNU}_S = 0.92 \\
U^{KNU}_{S,1} &= (0.92)^1 \cdot [0 - 2.5 \cdot 0.3] = 0.92 \cdot (-0.75) = -0.69
\end{align*}

\textbf{Component: $U^{IDN}_{E,0}$ (Identity Emotional, Instant)}
\begin{align*}
r^{IDN}_E &= 0.7 \quad \text{(``I'm a family person'')} \\
x^{IDN}_{E,0} &= 0.4 \quad \text{(guilt about leaving family)} \\
G^{IDN}_{E,0} &= 0 \\
P^{IDN}_{E,0} &= 0.7 - 0.4 = 0.3 \\
\lambda^{IDN}_E &= 2.8, \quad \delta^{IDN}_E = 1.0 \text{ (instant)} \\
U^{IDN}_{E,0} &= (1.0)^0 \cdot [0 - 2.8 \cdot 0.3] = -0.84
\end{align*}

\paragraph{Interpretation.}
The financial gain (positive INU) is offset by social pain (negative KNU)
and identity pain (negative IDN). The uniform calculation logic allows
direct comparison across categories and dimensions.

%%%%%%%%%%%%%%%%%%%%%%%%%%%%%%%%%%%%%%%%%%%%%%%%%%%%%%%%%%%%%%%%%%%%%%%%%%%%%%%
\subsubsection{Summary: Uniform Calculation Logic}
%%%%%%%%%%%%%%%%%%%%%%%%%%%%%%%%%%%%%%%%%%%%%%%%%%%%%%%%%%%%%%%%%%%%%%%%%%%%%%%

\begin{tcolorbox}[colback=gray!5!white, colframe=gray!75!black, title=Section 10.4 Summary]
\textbf{The Five-Step Calculation Procedure:}
\begin{enumerate}
\item Determine reference point $r^i_d$
\item Observe outcome $x^i_d$
\item Calculate Gain $G = \max(0, x - r)$ and Pain $P = \max(0, r - x)$
\item Apply loss aversion: $U_{raw} = G - \lambda \cdot P$
\item Apply discounting: $U = \delta^t \cdot U_{raw}$
\end{enumerate}

\textbf{The Master Formula:}
$$\boxed{U^i_{d,t} = (\delta^i_d)^t \cdot \left[ G^i_{d,t} - \lambda^i_d \cdot P^i_{d,t} \right]}$$

\textbf{Key Properties:}
\begin{itemize}
\item \textbf{Universal:} Same formula for all 144 components
\item \textbf{Category-specific parameters:} $r$, $\lambda$, $\delta$ vary by INU/KNU/IDN
\item \textbf{Dimension-specific parameters:} $\lambda$, $\delta$ vary by FEPSDE
\item \textbf{Context-modulated:} All parameters respond to $\Psi$
\end{itemize}

\textbf{The Symmetry Principle Confirmed:}
INU, KNU, and IDN use identical calculation logic.
The difference lies entirely in parameter values---especially reference point formation.
\end{tcolorbox}

\vspace{1em}
\hrule
\vspace{0.5em}
\noindent\textit{Next section: 10.5 Reference Point Formation}

%%%%%%%%%%%%%%%%%%%%%%%%%%%%%%%%%%%%%%%%%%%%%%%%%%%%%%%%%%%%%%%%%%%%%%%%%%%%%%%
